\statement{A.136}{Proved}{The second fixed axis of the actual two-strand reflections}
Retain the ORIGINAL negative \(\widehat{GT}_0\) two-strand action
of A.130 on \(E=\widehat F(X,Y)\). With the actual involutions
\(\beta,\theta\) of that theorem, put
\[
 \chi=\beta\theta,\qquad t_*=\operatorname{Ad}(X)\theta .
\]
Then
\[
 E^\chi=\overline{\langle X\rangle}\cong\widehat{\mathbb Z},
 \qquad H^1(C_2,E';\chi)=* .
 \tag{A.136.1}
\]
Together with A.130, this gives the two actual fixed axes
\[
 E^{\operatorname{Ad}(X)\beta}
       =\overline{\langle a_B\rangle},\quad [a_B]=[Y],
 \qquad E^{\beta\theta}=\overline{\langle X\rangle}.
 \tag{A.136.2}
\]
These equalities do not assert that \(X,a_B\) generate \(E\), or that
the specified canonical word \(a(k)\) equals \(a_B\).
\begin{proof}
\emph{1. The actual action and its cusp normalization.}
Put \(K=[Y,X]\), and use the exact A.130 words
\[
 \Phi=\phi(X,K),\qquad Q=j(\Phi),\qquad
 \theta=\operatorname{Ad}(Y^{-1})j,\qquad
 \beta=\operatorname{Ad}(\Phi^{-1})T .
\]
All commutators mean \(aba^{-1}b^{-1}\). The proved two-strand
construction gives
\[
\begin{aligned}
 \beta(X)&=\theta(X)=X^{-1},\\
 \beta(Y)&=\Phi^{-1}K^{-1}Y\theta(\Phi),&
 \beta(K)&=\Phi^{-1}K^{-1}\Phi,\\
 \theta(K)&=Y^{-1}KY,&
 \theta(\Phi)&=Y^{-1}QY.
\end{aligned}
\tag{A.136.3}
\]
The involutions \(\beta,\theta\) commute. Moreover \(t_*^2=1\)
and \(E^{t_*}=1\), by A.130's actual finite-factor free profinite
product centralizer argument. Thus \(\chi\) is an involution,
\(\chi(X)=X\), and its abelian action is
\(\operatorname{diag}(1,-1)\) in the ordered coordinates \(X,Y\).

Direct substitution, without changing the based cusp, gives
\[
\begin{aligned}
 \chi(K)
 &=\beta(Y)^{-1}\Phi^{-1}K^{-1}\Phi\beta(Y)\\
 &=\theta(\Phi)^{-1}Y^{-1}K^{-1}Y\theta(\Phi)\\
 &=(QY)^{-1}K^{-1}(QY).
\end{aligned}
\tag{A.136.4}
\]
Set \(g=QY\) and \(\sigma_*=\operatorname{Ad}(g)\chi\).
Then \(\sigma_*(K)=K^{-1}\). Initially its square is the inner
automorphism with coefficient \(\Delta=g\chi(g)\).
Squaring its action on \(K\) shows
\(\Delta\in C_E(K)=\overline{\langle K\rangle}\), by C.5.
Write \(\Delta=K^z\), \(z\in\widehat{\mathbb Z}\).
There is also the literal equality
\[
 \sigma_*(\Delta)
 =g\chi(g\chi(g))g^{-1}=g\chi(g)=\Delta.
\]
But \(\sigma_*(K^z)=K^{-z}\). The cusp is a faithfully
parametrized copy of \(\widehat{\mathbb Z}\), so \(2z=0\)
implies \(z=0\). Consequently \(\Delta=1\), and
\(\sigma_*\) is an actual involution.

In \(E\rtimes\langle\chi\rangle\), the section \(g\chi\)
is the cusp-normalized section. Since \(Q\in E'\), its coefficient
has homology \([Y]\). Relative to that section, \(\chi\) has
coefficient \(g^{-1}\), whose negative eigencoordinate is odd in
the \(Y\) direction. A coboundary has even negative
eigencoordinate. The two sections therefore have different
modulo-two homology labels. C.8 has exactly two section classes
at this actual full cusp and dyadic action; hence \(\chi\) is
the noncusp class. Its fixed homology line is the \(X\) line.
This identifies the branch without replacing the actual action
by a standard reflection.

\par\smallskip\noindent
\emph{2. Eliminate the complete pro-odd fixed kernel.}
The involutions \(\chi,t_*\) commute. Indeed \(\chi\) fixes
\(X\) and commutes with \(\theta\), whence
\[
 \chi t_*=\operatorname{Ad}(X)\chi\theta
          =\operatorname{Ad}(X)\beta=t_*\chi .
\]
Let \(B=E^\chi\). The actual primitive generator \(X\) belongs
to \(B\). The precise C.11--C.12 conclusions give
\[
\begin{gathered}
 \exp_X:B\twoheadrightarrow\widehat{\mathbb Z},\qquad
 M=\ker(\exp_X|_B)=B\cap E',\\
 B=M\rtimes\overline{\langle X\rangle},\qquad
 M\text{ is pro-odd}.
\end{gathered}
\tag{A.136.5}
\]
A fixed abelian vector has zero \(Y\) coordinate, which explains
the equality defining \(M\). The involution \(t_*\) preserves
\(B,M\), sends \(X\) to \(X^{-1}\), and has no fixed element in
\(M\).

Here is the full finite-to-compact step. Let \(J\) be an open
normal \(t_*\)-stable subgroup of \(M\). If \(xJ\) is fixed,
put \(n=x^{-1}t_*(x)\in J\). Then \(t_*(n)=n^{-1}\).
The closed procyclic group generated by \(n\) is pro-odd,
so its continuous half-power exists. For \(y=n^{-1/2}\in J\),
\[
 y^{-1}t_*(y)=n,\qquad t_*(xy^{-1})=xy^{-1}.
\]
The fixed subgroup of \(M\) is trivial, so \(xJ=J\).
Thus the finite group \(M/J\) has a fixed-point-free involution.
On this finite group the map \(h\mapsto h^{-1}t_*(h)\) is
injective, hence onto. The involution sends every such value to
its inverse. Therefore it is inversion on the whole group;
inversion being an automorphism forces that group to be abelian.
The subgroups \(J\) are cofinal, by intersecting each open normal
subgroup with its \(t_*\)-image. It follows that \(M\) is abelian
and \(t_*(m)=m^{-1}\) for every \(m\in M\).

Let \(\lambda=\operatorname{Ad}(X)|_M\). Applying \(t_*\) to
\(XmX^{-1}\) now gives
\[
 \lambda(m^{-1})=\lambda^{-1}(m^{-1}).
\]
Hence \(\lambda^2=1\), so \(X^2\) centralizes \(M\).
C.5, applied to the actual primitive \(X\) and its nonzero
integral power \(2\), gives
\(C_E(X^2)=\overline{\langle X\rangle}\).
Thus \(M\subseteq\overline{\langle X\rangle}\).
The character \(\exp_X\) is injective on that procyclic group
and vanishes on \(M\); therefore \(M=1\).
This proves the first equality in (A.136.1).

\par\smallskip\noindent
\emph{3. The complete derived cohomology and its scope.}
Use cocycles \(z\chi(z)=1\), with equivalence
\(z\sim h^{-1}z\chi(h)\).
An \(E'\)-valued cocycle has zero abelian coefficient, so its
section has the same modulo-two label as the noncusp section
\(\chi\). By C.8 it is an ambient coboundary
\(z=h^{-1}\chi(h)\). Abelianization forces the \(Y\) coordinate
of \(h\) to vanish, so \([h]=n[X]\).
Replace \(h\) by \(X^{-n}h\). This does not change its
coboundary, because \(X\) is \(\chi\)-fixed, and the new
conjugator belongs to \(E'\). Thus every derived cocycle is a
derived coboundary, proving the second assertion in (A.136.1).

The canonical family \(a(s)\) of A.129 satisfies
\(\chi(a(s))=a(s)^{-1}\), but has homology \([Y]\).
It lies in the other section branch, so the derived cohomology
conclusion cannot be applied to it. The remaining specified
boundary equality is still \(a(k)=a_B\).
\end{proof}

\paragraph{Source and unchanged decision scope.}
This proof consolidates the accepted entire source
\texttt{work\_DE/actual\_two\_strand\_second\_fixed\_axis37.md}
(SHA-256 beginning \texttt{de4444f4}), using the already proved
full scopes of A.130 and C.5, C.8, C.11, C.12.
It uses no three-strand action for arbitrary
\(\widehat{GT}_0\), no finite target search, and no generation
inference from the two fixed axes. Main A and the original CM
decision remain Open.
